\documentclass[11pt,a4paper]{article}
\usepackage[margin=23mm]{geometry}
\usepackage{amsmath,amssymb,amsthm}
\setlength{\emergencystretch}{2em}
\newtheorem{theorem}{Theorem}
\newcommand{\gen}[1]{\langle#1\rangle}
\newcommand{\Alt}{\operatorname{Alt}}
\title{An internal orbit with deficit two}
\author{Human Compression project: L044 continuation}
\date{}
\begin{document}\maketitle
Let $G=\gen{\rho_0,\ldots,\rho_{r-1}}\le\Alt(\Omega)$
be transitive with a connected string C-marking and all
maximal omissions intransitive. Assume $n=2r+1$ and
the actual transitive six-point head on $D=\{1,\ldots,6\}$,
with the literal interface
\[
 \tau=\rho_3=(6\ b)\delta,\quad
 R=G_{\ge3}=\Alt(\Omega\setminus\{1,\ldots,5\}),\quad
 \sigma=\rho_4,\quad \sigma|_D=1.
\]
The involution $\delta$ fixes $D,b$. Put $H=G_{\ge5}$
and $K=G_{>5}$, and suppose $H$ fixes exactly $D\cup\{b\}$.
All eligible edges below are eligible for the full ambient
omissions. Use actual normalized assignments of two such
edges of each tail colour to a single $H$-orbit.

\begin{theorem}
Suppose the two nontrivial $H$-orbits are $V$ of size
$2s$ and $O$ of size $2q$, where $s\ge2$ and $q\ge3$
is odd. Assume $\delta$ preserves $O$ and is free there,
that $K$ is intransitive on $O$, and that
\[
        5\in F_V,\qquad |F_V|=s,\qquad |F_O|=q-2.
\]
Then this configuration is impossible.
\end{theorem}
\begin{proof}
The acquired actual head-overlap argument gives
\[
 x=b^\sigma\ne b,\qquad \delta x=x,\qquad Kx=x.
                                                        \tag{1}
\]
These conclusions use natural $R$ and ambient IP, without
requiring $G_{\ge4}$ transitive. The point $x$ lies in $V$.
Since $\delta$ centralizes $H$ and fixes $x$, it fixes
all of $V$. The complete selected-colour equality makes
$G_{F_V}$ transitive on $V$. The acquired fixed-anchor
lemma then gives
\[
 F_V=\{5,\ldots,a-1\},\quad F_O=\{a,\ldots,r-1\},
 \quad a=s+5\ge7,\quad r=s+q+3.                       \tag{2}
\]
Set $T=G_{F_O}$ and $P=G_{\le a-1}$.

We first record two literal permutation facts needed
to justify projection. The prefix $P$ contains the pure
group $\Alt(V)$. Here is the degree-independent part of
the acquired prefix-kernel proof. Conjugate $\tau$ by
the transitive head and multiply two conjugates; the
exterior $\delta$ cancels, supplying the three-cycles
of the pure $\Alt(D\cup\{b\})$. Conjugation by $\sigma$
merges that support with $D\cup\{x\}$ and supplies
$\Alt(D\cup\{b,x\})$. Inductively a new mark fixes
the earlier support two indices back, because it
centralizes its natural alternating group. The two
current supports therefore overlap in at least seven
points. The complete support-merging lemma gives
the pure alternating group on the enlarged actual
head orbit. The head orbit of $P$ contains $x$ and
all its $G_{F_V}$-translates, hence all of $V$.
This proves the asserted inclusion. No even-degree
budget or rank induction enters these steps.

Secondly, every mark of $T$ except possibly $\rho_a$
fixes $V$ pointwise. It fixes $x$ by (1) and commutes
with the transitive group $G_{F_V}$ by (2). Thus all
$T$-orbits in $V$ have size at most two. All its marks
also fix $D,b$.

Project $O$ onto the $q$ pairs of $\delta$. Select one
eligible edge for each colour in $F_O$. Their projections
form a forest $\mathcal F$ with $q-2$ edges and two
components: an alternative projected path lifts using
$\tau$ and the other selected colours, contradicting
ambient eligibility. This is the complete pair-forest
argument, before any equality conclusion is used.
Each component lies inside a $T$-orbit on pairs.
Consequently that action has either one or two orbits.

The pair image of $\rho_5$ is trivial. Indeed $\rho_5$
commutes with $T$. If $T$ is transitive on pairs,
an involution of the odd set of pairs fixes a point,
and a centralizing permutation fixing one point fixes
them all. Otherwise its two orbits are the two forest
components; their sizes are unequal since $q$ is odd,
so $\rho_5$ preserves each. If it moved a vertex to
another vertex in the same component, join them by
the unique selected tree path. Any edge of that path,
of colour $j\ne5$, would acquire an alternative path
using $\rho_5$, the other tree edges and $\tau$, all
in $G_j$. This is again impossible.

Since $H$ is transitive on $O$, its pair action is
transitive, and the same now holds for $K$. Its assumed
intransitivity on points gives exactly two $K$-orbits
of size $q$, each meeting each pair once. Any element
of $K$ with trivial pair image fixes both rows
pointwise. In particular the kernel of the pair
projection on $T$ fixes $O$. It fixes $D,b$ as well,
and hence is a pure even permutation of $V$. The kernel
therefore lies in $\Alt(V)\le P$. Actual ambient IP
gives $T\cap P=1$, so the projection on $T$ is injective.
The projected consecutive marking thus retains both
IP and all its actual noncommutations.

Suppose first that $T$ is transitive on pairs. Its
point orbits in $O$ are then the two rows of size $q$.
The mark $\sigma$ centralizes $T$ and must carry a
$T$-orbit to one of the same size. Outside $O$ its
orbits have size at most two, whereas $q\ge3$.
Thus $\sigma$ preserves $O$. So do $\tau$ and $H$,
contrary to the natural transitivity of $R$.

Hence $T$ has exactly the two forest components as
pair orbits. Every projected mark is precisely its
selected transposition. To see this, an additional
edge lies within a forest component. If its tree
path contains an edge of a different colour, it
gives the same forbidden alternative path. A path
using only its own colour consists of the single
selected edge. There are no other possibilities.
The actual marking is connected and its projection
is faithful. Thus consecutive projected marks do
not commute, so all selected edges belong to one
component, and the other is a singleton. The tree
component is a path: three incident edges would
require three distinct indices to be pairwise
consecutive. Its successive indices are consecutive
in order. We have proved
\[
       T\text{ on pairs is the natural }S_{q-1}
       \text{ and one fixed pair}.                    \tag{3}
\]
On $O$, it has two orbits of size $q-1$ and two fixed
points, using the faithful rows of $K$.

The permutation $\sigma$ cannot preserve $O$, since
$\tau,H$ preserve it and generate $R$ together with
$\sigma$. As $\sigma$ centralizes $K$, a point
crossing out of $O$ brings an entire $K$-row of size
$q$ to a $K$-orbit outside $O$, necessarily in $V$.
On that row $T$ has an orbit of size $q-1$, by (3).
All its orbits in $V$ have size at most two.
It follows that $q-1\le2$, and therefore $q=3$.

Now $a=r-1$, $T=\gen t$ with $t=\rho_a$, and its
pair image is one transposition of three pairs.
The transitive pair image of $K$ contains that
transposition, and is therefore $S_3$. Put
$u=\rho_{a-1}$. For $6\le i\le a-2$, commutation
with $t$ shows that the pair image of $\rho_i$ is
either the identity or the image of $t$. In the
second case $\rho_i t$ fixes $O$, by faithfulness
of the literal $K$-action on its two rows. It is
then a pure even permutation of $V$, hence belongs
to $\Alt(V)\le P$. Since $\rho_i\in P$, this would
give $t\in P$, contradicting actual independence.
Consequently every such $\rho_i$ fixes $O$ pointwise.
The two remaining images, those of $u,t$, generate
$S_3$, so they are distinct transpositions.
The group $Q=\gen{u,t}$ thus has exactly two natural
three-point orbits on $O$. We make no assertion
that its whole action outside $O$ is faithful.

For $4\le l\le a-3$, put $Z_l=O^{G_{\le l}}$, the
union of the indicated translates. The prefix
centralizes $Q$, so $Z_l$ is a union of actual
$Q$-orbits of size three. The group $Q$ fixes
$D,b,x$, eight points disjoint from every $Z_l$.
Since $\sigma$ moves a $Q$-orbit out of $O$,
$|Z_4|\ge9$.
For $4\le l\le a-4$, every mark of index at least
$l+2$, other than $u,t$, fixes $O$ pointwise and
commutes with $G_{\le l}$; it therefore fixes $Z_l$
pointwise. The group $Q$ preserves $Z_l$.
If $\rho_{l+1}$ also preserved $Z_l$, then all of
$G$ would preserve that proper nonempty set.
Transitivity therefore forces enlargement, and
commutation with $Q$ makes the increment at least
one whole orbit of size three. We obtain
\[
 |Z_{a-3}|\ge9+3(a-7)=3a-12,\qquad
 n\ge(3a-12)+8=3a-4.                                  \tag{4}
\]
On the other hand $r=a+1$ gives $n=2a+3$.
Thus $a\le7$, and (2) forces $a=7$ and $s=2$.
But $F_V=\{5,6\}$ now assigns two distinct
$6$-edges to the four-point set $V$. An involution
with those edges fixes no point of $V$, contrary
to $\rho_6x=x$. This final contradiction proves
the theorem.
\end{proof}

\paragraph{Attribution and remaining scope.}
The actual head interface, selected-colour equality,
fixed-anchor lemma, support-merging kernel proof and
pair-forest argument are acquired inputs. The new
steps are the faithful terminal projection through
the actual pure kernel, the two-component analysis,
and the final three-point orbit growth. Root's earlier
support-kernel and orbit-matching contributions are
retained in the acquired material. There is no new
computation, rank induction or formal premise.
This excludes the stated normalized budget only;
split-orbit assignments, global attachment and the
complete L044 recovery remain separate. No original,
question, formal or numerical credit is claimed.
\end{document}
